\documentclass{article}

% Language setting
\usepackage[english]{babel}

% Set page size and margins
\usepackage[letterpaper,top=2cm,bottom=2cm,left=3cm,right=3cm,marginparwidth=1.75cm]{geometry}

% Useful packages
\usepackage{amsmath}
\usepackage{graphicx}
\usepackage[colorlinks=true, allcolors=blue]{hyperref}

\title{LE 105: Assignment 3}
\author{Sreyash Gupta}

\begin{document}

\maketitle

\section*{Question 1}
\subsection*{Part a}
A solar eclipse is a fascinating celestial event that occurs when the Moon moves in such a way that it positions itself directly between the Earth and the Sun. In this special alignment, the Moon blocks the light coming from the Sun and prevents it from reaching certain parts of the Earth’s surface. As a result, observers on Earth can see either a partial or complete covering of the Sun by the Moon. This precise positioning requires the Sun, Moon, and Earth to be collinear, that is, lying along the same straight line.  

In contrast, a lunar eclipse takes place when the Earth comes between the Sun and the Moon. Since the Moon does not emit any light of its own, it only shines by reflecting the sunlight that falls upon it. During a lunar eclipse, the Earth obstructs this sunlight, casting its shadow onto the Moon. Consequently, the Moon appears darkened or even completely invisible for a certain period of time. Just as in the case of a solar eclipse, this alignment also requires the Sun, Earth, and Moon to be exactly collinear.

\subsection*{Part b}
The occurrence of eclipses is also tied to specific phases of the lunar cycle. A solar eclipse always happens on the new moon day, known as \textit{Amavasya}, when the Moon lies between the Earth and the Sun. On the other hand, a lunar eclipse can only occur on a full moon day, referred to as \textit{Purnima}, when the Earth positions itself between the Sun and the Moon. Thus, the lunar phases play a crucial role in determining when eclipses are possible.

\subsection*{Part c}
One might wonder why, despite the regular cycle of new moons and full moons, eclipses do not occur every month. The primary reason lies in the fact that the Moon’s orbital plane is not perfectly aligned with the Earth’s orbital plane around the Sun. The orbit of the Moon is inclined by approximately 5 degrees relative to the Earth’s orbit (the ecliptic plane). Because of this tilt, the Moon is usually either slightly above or slightly below the imaginary straight line connecting the Sun and the Earth. Only when the Moon happens to cross this line at points known as the nodes do eclipses occur. Hence, eclipses are relatively rare events, occurring only a few times each year instead of every month.

\subsection*{Part d}
The nature and duration of eclipses are largely determined by the relative sizes and distances of the Sun, Moon, and Earth. Although the Sun is about 400 times larger in diameter than the Moon, it is also about 400 times farther from the Earth. This remarkable coincidence makes the Sun and Moon appear to be almost the same size when viewed from Earth’s surface. It is precisely this similarity in apparent size that makes solar eclipses possible.  

In the case of a lunar eclipse, the Earth is approximately 3.67 times larger in diameter than the Moon. As a result, the Earth’s shadow is wide enough to completely cover the Moon during a total lunar eclipse. This explains why a lunar eclipse typically lasts longer than a solar eclipse. By contrast, during a solar eclipse, the smaller size of the Moon means that its shadow covers only a narrow region of the Earth, making solar eclipses much shorter in duration. These differences highlight the delicate balance of cosmic geometry that allows us to witness such awe-inspiring events.

\section*{Question 2}
The given relation can be simplified as follows:
\[
354.3672n + 29.5306m = 365.2594n
\]
\[
\Rightarrow 10.8922n = 29.5306m
\]
\[
\Rightarrow \frac{n}{m} = \frac{29.5306}{10.8922}
\]

This ratio indicates that for approximately every 30 lunar months, there are about 11 solar months. In other words, the passage of time as measured by the Moon (lunar months) and as measured by the Sun (solar years) do not align perfectly.  

Over the span of 19 solar years, the difference becomes even more pronounced. During this period, there are about 228 solar months, but nearly 235 lunar months. This mismatch of 7 months means that a purely lunar calendar would gradually fall out of sync with the solar calendar, which is tied to the seasons. To correct for this discrepancy and to keep the lunar calendar in harmony with the solar year, an intercalary or leap month is periodically added. Specifically, 7 leap months are introduced over a 19-year cycle. This system ensures that festivals and agricultural seasons, which are dependent on the solar year, continue to occur at the appropriate times when using a lunar-based calendar.

\end{document}
